\documentclass[12pt]{article}
\usepackage[utf8]{inputenc}
\usepackage[T1]{fontenc}
\usepackage{hyperref}
\title{W\'ara Language Resources for K-12 Curriculum Development}
\author{Compiled by Claude Sonnet 5 (claude-sonnet-5) for Andr\'as Kornai}
\date{}
\begin{document}
\maketitle

\section{Introduction}
W\'ara (ISO 639-3 \texttt{tci}), also called Upper Morehead, Vara, Wara or Rouku, is a Papuan language of the Yam family in Western Province, Papua New Guinea. Its varieties include K\'omnjo (Rouku), Anta and W\`er\'e, and some are sometimes treated as separate languages. Wikipedia, citing Ethnologue 2018, gives about 350 speakers of W\'ara proper and about 930 for the wider group. Do not confuse it with Warapu (\texttt{wra}) or with Awara. This survey found very little material, and nothing designed for schools.

\section{Primary Educational Resources}
No primers, textbooks or school programmes turned up in the searches.

\section{Linguistic and Reference Materials}
\begin{itemize}
\item Wikipedia, Upper Morehead language: \url{https://en.wikipedia.org/wiki/Upper_Morehead_language}
\item Wara Komnjo numerals page (senary yam counting system): \url{https://lingweb.eva.mpg.de/channumerals/Wara-Komnjo.htm}
\item Glottolog reference: Wurm and Sarsa, ``Organised phonology: W\'ara language, Western province'': \url{https://glottolog.org/resource/reference/id/313393}
\item ASJP, Wara (PNG) word list: \url{https://asjp.clld.org/languages/WARA_PNG}
\end{itemize}

\section{Religious and Community Resources}
None found.

\section{Digital and Audio Resources}
\begin{itemize}
\item Global Recordings Network, language entry for \texttt{tci}. I did not confirm what recordings it holds: \url{https://globalrecordings.net/en/language/tci}
\item Endangered Languages Project entry (I did not confirm that this page concerns W\'ara): \url{https://endangeredlanguages.com/lang/10744?hl=de}
\end{itemize}

\section{Curriculum Development Context}
Little published material exists. The yam-based numeral system could be a culturally grounded starting point for mathematics. A usable orthography and teaching materials would have to be developed with the community.
\end{document}